\section{Relation to the original work}\label{sec:origins}

\paragraph{Source-to-structure networks.} \citet{bromberg2020hierarchy} asked how hierarchy arises from local adaptation
to where information appears. Its premise --- that the location of information shapes who listens to whom, producing
hierarchy --- holds, within a window of commitment $\kappa<T<\Thalf$. Its hard selection cut imposes hierarchy even when
the signal carries no information (with a random source the cut still climbs toward the same region of the morphospace),
and its hierarchy measure needed two corrections (\cref{app:errata}).

\paragraph{Optimal learning in noisy niches.} \citet{lee2024constructing} predicted that optimal memory grows like
$\tauE^{1/2}$ and that agents who stabilize their environment create a shared good that can favour monopolies. The
attention network reproduces the scaling once signals are noisy (exponent 0.43); with memory and pace coupled, it settles
at a stable point with free-riding; but paying for attention does not entrench hubs. Their agents correspond to the
evidence rule of \cref{sec:learning}; the network adds relays, which help under noise.

\paragraph{Adaptive transport networks.} Fluctuating loads keep loops
\citep{corson2010fluctuations,katifori2010damage}. In the attention model this holds only for a flow-based reward; with
closeness, a fluctuating source merely delays the transition to some tree. The kind of budget, not the reward, decides
robustness and cost, and transport dissipation prices the structure that the attention model builds without paying for it.

\paragraph{The synthesis note's predictions.} Three regimes (distributed, hierarchy, over-betting): confirmed, with exact
boundaries. Memory scaling: confirmed with noise. Paid-attention monopolies: not found. The update interval as a sampling
window: holds only when rewards are averaged rather than distances summed.

\paragraph{Related learning dynamics.} The update belongs to a family studied in game theory and statistical physics:
replicator dynamics with memory loss \citep{sato2003coupled}, experience-weighted attraction
\citep{camerer1999experience}, reinforcement-driven network formation
\citep{skyrms2000dynamic,bala2000noncooperative}, and adaptive networks more broadly \citep{gross2008adaptive}.
